% --------------------
% Modelo Oficial para Projeto de Qualificação PPGEA 
% 
% Francisco Reinaldo 
%		(https://orcid.org/0000-0001-6161-6755)
%		(http://lattes.cnpq.br/7401534350061823)
% --------------------
% 29/06/20 finalização com texto explicativo incorporado
% --------------------
% Agradecimentos a Overleaf pela oportunidade 
% --------------------
% Mestrado em Engenharia Ambiental -  Francisco Beltrão/PR
% --------------------
\documentclass[oneside, 12pt,article, a4paper]{memoir}

%---- ---- Fontes
\usepackage[T1]{fontenc}
\usepackage[utf8]{inputenc}
\usepackage{lmodern,microtype}
\usepackage{fix-cm}
\usepackage[brazil]{babel}


%---- ---- Margens
\setlrmarginsandblock{3.5cm}{2.5cm}{*}
\setulmarginsandblock{3.5cm}{2.5cm}{*}
\checkandfixthelayout 


%---- ---- Texto
%Alinhamento
\usepackage{ragged2e} %\justifying


%---- ---- Figuras
\usepackage{graphicx}


%---- ---- Tabelas
%tabelas com cor de fundo 
\usepackage[svgnames,dvipsnames,x11names,table,xcdraw]{xcolor} 


%---- ---- ToC
%Centralizar
\renewcommand{\printtoctitle}{\centering\bfseries}

%Renomear
\addto\captionsbrazil{%
       %\renewcommand\listfigurename{\bfseries Lista de Ilustrações}
       %\renewcommand\listtablename{\bfseries Lista de tabelas}
       \renewcommand\contentsname{\bfseries SUMÁRIO\vskip2em}
       \renewcommand{\bibname}{\bfseries REFER\^ENCIAS}
       %\renewcommand{\indexname}{\uppercase{\bfseries \'Indice}
}


%---- ---- Parágrafo
%Novo ambiente
\usepackage{quoting}
\newenvironment{fdrosto}{\begin{quoting}[rightmargin=0cm,leftmargin=7cm]\small\noindent\hspace{-0.43em}\justifying \setlength\parindent{0pt}}{\end{quoting}}

%Espaçamento 1.5
\OnehalfSpacing

\usepackage{indentfirst}


%---- ---- Estilo bibliografico
\usepackage[alf,versalete]{abntex2cite}
%\bibliographystyle{unsrt}


%---- ---- Memoir: Troca fontsize 
%\chapterstyle{section}
\renewcommand*{\chapnumfont}{\normalfont\large\bfseries}
\renewcommand*{\chaptitlefont}{\MakeUppercase{\normalfont\large\bfseries}}

%in section
\newcommand*\lowercasecapitals[1]{\large\bfseries\scshape#1}
\setsecheadstyle{\lowercasecapitals}


%---- ---- Macros
\def\mestrando{NOME DO MESTRANDO}
\def\tituloprojeto{TÍTULO DO PROJETO}
\def\orientacao{Prof. Doutor em Eng. Francisco Antonio Fernandes Reinaldo}
\def\coorientacao{Profa. Doutora Sheila Oro}
\def\linhapesquisa{Análise Ambiental ou Tecnologia Ambiental}
%-------------------------------------

\begin{document}
\pagestyle{empty}

%---------CAPA 
\begin{center} 
    UNIVERSIDADE TECNOLÓGICA FEDERAL DO PARANÁ
    
    CAMPUS FRANCISCO BELTRÃO
    
    PROGRAMA DE PÓS-GRADUAÇÃO EM ENGENHARIA AMBIENTAL: 
    
    ANÁLISE E TECNOLOGIA AMBIENTAL
    
    \vskip15em \mestrando
    
    \begin{vplace}[0.7] \tituloprojeto \end{vplace}
    
    \vfill
    FRANCISCO BELTRÃO/PR \\ \the\year
\end{center}
%---------CAPA

%---------ROSTO
\newpage
\begin{center}
\mestrando \vskip 8em \tituloprojeto \vskip 8em

\begin{fdrosto}
Projeto de dissertação apresentado ao Programa de Pós-Graduação em Engenharia Ambiental: Análise e Tecnologia Ambiental como requisito parcial à qualificação do curso de mestrado. \\

Linha de Pesquisa: \linhapesquisa

Orientador: \orientacao

Co-orientadora: \coorientacao
\end{fdrosto}

\vfill Francisco Beltrão/PR \\ \the\year 
\end{center}
%---------ROSTO

%---------SUMARIO
\newpage
\begin{KeepFromToc}
  \tableofcontents
  \thispagestyle{empty}
\end{KeepFromToc}
%---------SUMARIO

%---------CORPUS
\newpage
\pagestyle{simple}


\section*{RESUMO} \setlength\parindent{0pt}

Deve-se fazer um resumo do projeto com um mínimo de 1000 e máximo de 2000 caracteres (com espaços). O texto é corrido, sem parágrafos e sem recuo. Texto com fonte Arial ou Times New Roman, tamanho 12, espaço 1,5. O projeto deve conter no máximo 20 páginas. 

{\vskip1em\bfseries PALAVRAS-CHAVE:} 
%(incluir cinco palavras-chave, que não constem no título do projeto)
Qualificação.
Mestrado. 
UTFPR.
Francisco Beltrão.
Pesquisa.


\section*{ABSTRACT}

Aqui deve-se apresentar o resumo traduzido para o inglês  com um mínimo de 1000 e máximo de 2000 caracteres (com espaços). 

O texto é corrido, sem parágrafos e sem recuo. Texto com fonte Arial ou Times New Roman, tamanho 12, espaço 1,5. O projeto deve conter no máximo 20 páginas. 

KEYWORDS: 
{\vskip1em\bfseries KEYWORDS:} 
%(incluir cinco palavras-chave, que não constem no título do projeto)
Qualify.
Master Degree. 
UTFPR.
Francisco Beltrão.
Pesquisa.

\vskip2em
\setlength\parindent{1.25cm}
\chapter{INTRODUÇÃO}

\begin{figure}[htbp]
\begin{center}
\caption{título}
\includegraphics[width=0.7\textwidth]{overleaf}\\
\label{etiqueta}
Fonte: Autoria própria.
\end{center}
\end{figure}


\section{Teste}


Para que os membros da banca\cite{atalhomodelo} possam entender o seu projeto, conforme Figura\ref{etiqueta}, apresente aqui o assunto que irá abordar em seu projeto, contextualizando o problema de pesquisa \cite[p.~150]{atalhomodelo} perante a literatura científica baseada nas principais referências bibliográficas da sua área. 

Uma boa introdução pode conter entre 2000 a 4000 caracteres (com espaços).

\chapter{OBJETIVOS}

Coloque aqui de forma explícita\cite{Sommerville2011} todos os objetivos de seu projeto (geral e específicos). 

O objetivo geral deve abranger os objetivos específicos que serão definidos neste item. Cada objetivo começa geralmente com um verbo no infinitivo.

\chapter{JUSTIFICATIVA}

Revele a importância e abrangência do assunto e da sua proposta de pesquisa. O problema abordado por sua pesquisa deve ser relevante o suficiente para merecer uma investigação científica.

\chapter{REVISÃO DA LITERATURA}

Neste item abordam-se os fundamentos teóricos da pesquisa. A revisão deve ser obtida a partir das referências bibliográficas, principalmente em periódicos científicos especializados, nacionais e internacionais. Evitar a utilização de literatura cinza (relatórios, trabalhos de conclusão de curso, anais de eventos, e similares). Evitar utilizar apud (citado por) porque não é recomendado para trabalhos de pós-graduação stricto sensu.

\chapter{MATERIAIS E MÉTODOS}

Neste item deve-se colocar todo o material e os métodos (referenciados) para obtenção dos resultados. Quando pertinente, um delineamento experimental (com etapas) deverá ser desenvolvido de forma a ficar bem estabelecido qual o direcionamento do projeto de pesquisa e resultados esperados. Deve-se indicar na metodologia como será o tratamento estatístico de seus dados – quais análises e testes estatísticos serão feitos. 
A metodologia deve ser clara de modo que o leitor, membro da banca, consiga entender e supostamente consiga desenvolver o projeto seguindo os passos indicados na metodologia.

\chapter{RESULTADOS PARCIAIS E DISCUSSÃO}

Aqui devem ser inseridos e discutidos os primeiros resultados já obtidos pelo aluno até o momento da qualificação. É permitido o uso de figuras, tabelas, gráficos, que sirvam para ilustrar e explicar da melhor forma possível os resultados já alcançados, mesmo que preliminares.
Após a apresentação dos resultados, estes precisam ser discutidos tendo como base a literatura científica nacional e internacional.

\chapter{PERSPECTIVAS FUTURAS}

Explique aqui quais os experimentos/etapas que ainda faltam ser realizados e qual sua previsão de concluir esses experimentos/etapas. 

\chapter{CRONOGRAMA}
Indique no cronograma de 24 meses as etapas que já foram ou serão executadas. As etapas podem ser adaptadas conforme a necessidade.


\begin{table}[!htb]
\footnotesize
%\resizebox{\linewidth}{!}{
\begin{tabular}{|p{4cm}|c|c|c|c|c|c|c|c|c|c|}
\hline
\textbf{2017}              & \multicolumn{1}{l|}{\textbf{Mar}} & \multicolumn{1}{l|}{\textbf{Abr}} & \multicolumn{1}{l|}{\textbf{Mai}} & \multicolumn{1}{l|}{\textbf{Jun}} & \multicolumn{1}{l|}{\textbf{Jul}} & \multicolumn{1}{l|}{\textbf{Ago}} & \multicolumn{1}{l|}{\textbf{Set}} & \multicolumn{1}{l|}{\textbf{Out}} & \multicolumn{1}{l|}{\textbf{Nov}} & \multicolumn{1}{l|}{\textbf{Dez}} \\ \hline
Elaboração do projeto      &                                   &                                   &                                   &                                   &                                   &                                   &                                   &                                   &                                   &                                   \\ \hline
Levantamento bibliográfico &                                   &                                   &                                   &                                   &                                   &                                   &                                   &                                   &                                   &                                   \\ \hline
Coleta de dados            &                                   &                                   &                                   &                                   &                                   &                                   &                                   &                                   &                                   &                                   \\ \hline
Análise de dados           &                                   &                                   &                                   &                                   &                                   &                                   &                                   &                                   &                                   &                                   \\ \hline
\end{tabular}
%}
\end{table}


%----------------------
\begin{table}[!htb]
\footnotesize
\centering
%\caption{tebela}
%\label{tab:my-table}
\begin{tabular}{|p{4cm}|c|c|c|c|c|c|c|c|c|c|}
\hline
\textbf{2018} &
  \multicolumn{1}{l|}{\textbf{Jan}} &
  \multicolumn{1}{l|}{\textbf{Fev}} &
  \multicolumn{1}{l|}{\textbf{Mar}} &
  \multicolumn{1}{l|}{\textbf{Abr}} &
  \multicolumn{1}{l|}{\textbf{Mai}} &
  \multicolumn{1}{l|}{\textbf{Jun}} &
  \multicolumn{1}{l|}{\textbf{Jul}} &
  \multicolumn{1}{l|}{\textbf{Ago}} &
  \multicolumn{1}{l|}{\textbf{Set}} &
  \multicolumn{1}{l|}{\textbf{Out}} \\ \hline
Levantamento bibliográfico                                                   &  &  &  &  &  &  &  &  &  &  \\ \hline
Coleta de dados                                                              &  &  &  &  &  &  &  &  &  &  \\ \hline
Análise de dados                                                             &  &  &  &  &  &  &  &  &  &  \\ \hline
Qualificação                                                                 &  &  &  &  &  &  &  &  &  &  \\ \hline
Participação em congresso para apresentação de resultados parciais ou finais &  &  &  &  &  &  &  &  &  &  \\ \hline
Redação da dissertação                                                       &  &  &  &  &  &  &  &  &  &  \\ \hline
Preparação e envio de artigo científico                                      &  &  &  &  &  &  &  &  &  &  \\ \hline
\end{tabular}
\end{table}


%------------------------
\begin{table}[!htb]
\footnotesize
%\resizebox{\linewidth}{!}{
%\fontsize{10}{12}\selectfont
\begin{tabular}{|p{4cm}|c|c|}
\hline
\textbf{2019}                                     & \multicolumn{1}{l|}{\textbf{Jan}} & \multicolumn{1}{l|}{\textbf{Fev}} \\ \hline
Defesa da dissertação                             &                                   &                                   \\ \hline
Correções após defesa para recebimento do diploma &                                   &                                   \\ \hline
\end{tabular}
%}
\end{table}

%---------CORPUS


\bibliography{acervo}

\end{document}